\documentclass[12pt]{article}
\usepackage[utf8]{inputenc}
\usepackage[T1]{fontenc}
\usepackage{hyperref}
\title{Tumak Language Resources for K-12 Curriculum Development}
\author{Compiled by Claude Sonnet 5 (claude-sonnet-5) for Andr\'as Kornai}
\date{}
\begin{document}
\maketitle

\section{Introduction}
Tumak (ISO 639-3 \texttt{tmc}, Glottocode tuma1260) is a Chadic (Afro-Asiatic) language of southern Chad, with about 25,000 speakers (1993 census figure on Wikipedia). Motun and Mawer are listed as dialects, though Blench treats them as separate languages. Sources place the speakers in Moyen-Chari/Koumra and Tandjil\'e/Mandoul. Most Motun speakers also use Sara.

\section{Primary Educational Resources}
No primers, textbooks, dictionaries or school programmes were found.

\section{Linguistic and Reference Materials}
\begin{itemize}
\item Wikipedia article on the Tumak language: \url{https://en.wikipedia.org/wiki/Tumak_language}
\item Glottolog map page for Tumak: \url{https://glottolog.org/resource/languoid/id/tuma1260.bigmap.html}
\item Britannica entry on the Tumak people: \url{https://www.britannica.com/topic/Tumak}
\end{itemize}

\section{Religious and Community Resources}
\begin{itemize}
\item Scripture Earth page for Tumak (Gospel of Mark on Bible.is, Global Recordings audio): \url{https://www.scriptureearth.org/00eng.php?iso=tmc}
\item Joshua Project entry for Tumak in Chad: \url{https://joshuaproject.net/people_groups/15630/CD}
\end{itemize}

\section{Digital and Audio Resources}
Scripture Earth lists audio recordings from Global Recordings Network (see above). I did not check the audio itself.

\section{Curriculum Development Context}
Little educational or linguistic documentation was found. Curriculum work would need new orthography and materials development, probably in cooperation with Sara-language education in the area.
\end{document}